% Fuentes: LOGBOOK.md EXP-002..EXP-010, DEC-006.
\section{Decisiones y resultados por fase}
Cada experimento cambia un solo factor respecto al ganador de la fase anterior, y se elige por el \emph{mean recall}
de validación. Medido con varias semillas, el ruido entre ejecuciones es de unos \res{seed-noise-pp}~pp; por eso las
últimas decisiones se toman con la media de \res{final-seeds} semillas, y las diferencias menores que ese ruido no se
consideran concluyentes~\cite{bouthillier}. El Cuadro~\ref{tab:phases} resume lo que se mantuvo en cada fase.

\begin{table}[h]
  \centering\small
  \safetable{generated/tables/phase_summary}
  \caption{Configuración que se mantiene al final de cada fase (validación; media $\pm$ desviación típica cuando hay
  varias semillas).}
  \label{tab:phases}
\end{table}

\subsection{Normalización y punto de parada (fase 1)}
\decisionpoint{con píxeles en $[0, 255]$ los logits son enormes y la softmax se satura~\cite{lecun}; y unos callbacks
que no actúan dejan sin usar la información de validación.}
{la pérdida de entrenamiento del baseline se queda en torno a \res{baseline-train-loss} durante todo el entrenamiento
(Figura~\ref{fig:normalisation}), y la \emph{learning rate} no cambia nunca.}
{escalar la entrada a $[0, 1]$ con una capa \texttt{Rescaling} baja la pérdida dos órdenes de magnitud
(Figura~\ref{fig:normalisation}) y sube el \emph{mean recall} \res{gain-normalisation}~pp.
Con paciencias acordes a la duración del entrenamiento, \emph{ReduceLROnPlateau} reduce la \emph{learning rate} al
estancarse la validación, desaparecen las oscilaciones y ninguna clase queda sin predecir (recall mínimo
\res{phase1-min-recall}). Vigilar el \emph{mean recall} dio el mejor resultado (\res{phase1-mean-recall}, frente a
\res{monitor-loss-mean-recall} vigilando la pérdida y \res{monitor-acc-mean-recall} vigilando la \emph{accuracy};
una sola semilla, diferencia pequeña), y es además la métrica de decisión.}
{\texttt{Rescaling} dentro del modelo; \emph{EarlyStopping} y \emph{ReduceLROnPlateau} sobre el \emph{mean recall}
de validación, y se conserva la mejor época.}

\textbf{Punto de parada.} En el modelo final (Figura~\ref{fig:stopping}, semilla~1) se ve cómo decide la validación: hasta la época
\res{stop-last-joint-epoch} train y validación suben juntas; al estancarse la validación, \emph{ReduceLROnPlateau}
divide la \emph{learning rate} por 10 y llega la mejor época (\res{stop-best-epoch}); después solo mejora train
(\emph{accuracy} \res{stop-train-acc-best} $\rightarrow$ \res{stop-train-acc-end}, con validación en
\res{stop-val-acc-best} $\rightarrow$ \res{stop-val-acc-end}), así que \emph{EarlyStopping} (paciencia
\res{final-es-patience}) para en la época \res{stop-epoch} y se conserva la \res{stop-best-epoch}.

\begin{figure}[!htbp]
  \centering
  \begin{minipage}[t]{0.485\linewidth}
    \centering
    \safegraphics[width=\linewidth]{figures/normalisation.png}{Pérdida con y sin escalado}
    \caption{\figcaption{normalisation}}
    \label{fig:normalisation}
  \end{minipage}\hfill
  \begin{minipage}[t]{0.485\linewidth}
    \centering
    \safegraphics[width=\linewidth]{figures/stopping.png}{Punto de parada}
    \caption{\figcaption{stopping}}
    \label{fig:stopping}
  \end{minipage}
\end{figure}

\subsection{Resolución de entrada (fase 2)}
\decisionpoint{a \res{res224-res}~$\times$~\res{res224-res}, una capa oculta (\res{res224-hidden}) tiene
\res{res224-params} parámetros para \res{n-train} imágenes: demasiada varianza.}
{al subir la resolución baja también el \emph{mean recall} de train (Figura~\ref{fig:resolution}) y la \emph{accuracy}
de train (\res{res224-train-acc} a \res{res224-res} píxeles):
es un problema de optimización, porque la mayoría de las ReLU queda inactiva.}
{el máximo está en \res{phase2-res}~$\times$~\res{phase2-res}, con un \emph{mean recall} de \res{phase2-mean-recall}
(+\res{gain-resolution}~pp sobre la fase~1), y se mantiene aunque cada resolución use su mejor \emph{learning rate}
(\res{res32-lr3e-4-mean-recall} a \res{res32-lr3e-4-res} y \res{res64-lr1e-4-mean-recall} a \res{res64-lr1e-4-res};
Figura~\ref{fig:resolution}).}
{entrada de \res{phase2-res}~$\times$~\res{phase2-res}, el mayor salto individual.}

\subsection{Arquitectura y activación (fases 3 y 4)}
\decisionpoint{más capacidad (más anchura o más capas) debería mejorar la validación mientras no haya sobreajuste.}
{con más parámetros la validación se estanca y train sigue subiendo (Figura~\ref{fig:capacity}): la red de
\res{big3layer-n-hidden} capas (\res{big3layer-hidden}) llega a una \emph{accuracy} de train de
\res{big3layer-train-acc} con el mismo \emph{mean recall} (\res{big3layer-mean-recall}) que \res{phase3-hidden}.}
{\res{phase3-hidden} (\res{phase3-mean-recall}, +\res{gain-depth}~pp sobre una sola capa) empata con la red mayor con
muchos menos parámetros (\res{phase3-params} frente a \res{big3layer-params}), y con el mismo presupuesto supera a una
sola capa ancha (\res{onelayer512-hidden}: \res{onelayer512-mean-recall}). LeakyReLU eliminó las ReLU muertas y llegó a
\res{phase4-mean-recall}.}
{\res{phase3-hidden} con LeakyReLU e inicialización de He.}

\begin{figure}[!htbp]
  \centering
  \begin{minipage}[t]{0.485\linewidth}
    \centering
    \safegraphics[width=\linewidth]{figures/resolution_sweep.png}{Barrido de resolución}
    \caption{\figcaption{resolution_sweep}}
    \label{fig:resolution}
  \end{minipage}\hfill
  \begin{minipage}[t]{0.485\linewidth}
    \centering
    \safegraphics[width=\linewidth]{figures/capacity.png}{Curva de capacidad}
    \caption{\figcaption{capacity}}
    \label{fig:capacity}
  \end{minipage}
\end{figure}

\subsection{Algoritmo de optimización y parámetros de entrenamiento (fase 5)}
\decisionpoint{la \emph{learning rate} es el hiperparámetro más influyente, y el optimizador y el batch fijan el coste
de cada experimento.}
{con hasta \res{adam-lr3e-4-epochs-cap} épocas, Adam es casi insensible a la \emph{learning rate}
(mejor \res{adam-lr3e-4-mean-recall}) y SGD con momentum tiene un óptimo marcado (\res{sgd-m09-mean-recall} a
\res{sgd-m09-lr}; Figura~\ref{fig:lr}). Repetidos con \res{final-seeds} semillas, los dos finalistas
empatan dentro del ruido.}
{Adam con $\beta_1$ = \res{final-beta1} y un batch de \res{final-batch} (\emph{learning rate} \res{final-lr}, escalada
con la raíz del batch) hace cada época \res{speedup-batch256} veces más rápida que con batch \res{phase4-batch}, a
cambio de unos \res{cost-batch256}~pp. Ningún \emph{schedule} ni otro ajuste del optimizador mejoró a
\emph{ReduceLROnPlateau}.}
{Adam con AMSGrad, $\beta_1$ = \res{final-beta1}, batch \res{final-batch}: la menor varianza entre semillas y la
rapidez permiten decidir con varias semillas.}

\begin{figure}[!htbp]
  \centering
  \safegraphics[width=0.5\linewidth]{figures/lr_sweep.png}{Barrido de learning rate}
  \caption{\figcaption{lr_sweep}}
  \label{fig:lr}
\end{figure}

